% =====================================================================
% Pricing & Risk Pipeline for Synthetic CDOs — Working paper
% Auteur : TOM (M2 Sciences Actuarielles, ISFA — Lyon 1, 2025-2026)
% Date   : 26 avril 2026
% =====================================================================
\documentclass[11pt,a4paper]{article}

% --- Encodage et langue ----------------------------------------------
\usepackage[utf8]{inputenc}
\usepackage[T1]{fontenc}
% French babel (fall back gracefully if texlive-lang-french missing)
\IfFileExists{french.ldf}{%
  \usepackage[french]{babel}%
  \@ifpackagelater{babel}{2017/05/06}{\frenchsetup{SmallCapsFigTabCaptions=false}}{}%
}{%
  \usepackage[english]{babel}%
  \PackageInfo{main}{French babel not installed; falling back to English layout.}%
}
\usepackage{lmodern}

% --- Mise en page ----------------------------------------------------
\usepackage[a4paper,margin=2.4cm,top=2.6cm,bottom=2.6cm]{geometry}
\usepackage{setspace}
\onehalfspacing
\usepackage{microtype}
\usepackage{parskip}

% --- Mathématiques ---------------------------------------------------
\usepackage{amsmath,amssymb,amsthm,mathtools}
\usepackage{bm}
% siunitx (optional — fall back to plain text if unavailable)
\makeatletter
\IfFileExists{siunitx.sty}{%
  \usepackage{siunitx}%
  \sisetup{output-decimal-marker={,},group-separator={\,}}%
}{%
  \long\def\SI##1##2{##1\,##2}%
  \long\def\si##1{\,##1}%
}
\makeatother

% --- Tableaux & figures ----------------------------------------------
\usepackage{graphicx}
\graphicspath{{images/}}
\usepackage{booktabs}
\usepackage{array}
\usepackage{multirow}
\usepackage{caption}
\captionsetup{font=small,labelfont=bf,format=plain,labelsep=period,
              justification=justified,singlelinecheck=false}
\usepackage{subcaption}
\usepackage{float}
\usepackage{longtable}
\usepackage{tabularx}

% --- Couleurs et boîtes ----------------------------------------------
\usepackage{xcolor}
\definecolor{cprimary}{HTML}{0B3B6F}
\definecolor{caccent}{HTML}{0C8A64}
\definecolor{cwarn}{HTML}{E65100}
\definecolor{cdanger}{HTML}{C62828}
\definecolor{csoft}{HTML}{F5F7FB}
\definecolor{cline}{HTML}{E1E6EF}

\usepackage[most]{tcolorbox}
\newtcolorbox{callout}[1][]{
  colback=csoft,colframe=caccent,
  boxrule=0pt,leftrule=3pt,arc=2pt,
  fontupper=\small,
  enhanced,breakable,
  title=#1,
  coltitle=cprimary,fonttitle=\bfseries\small,
}
\newtcolorbox{calloutwarn}[1][]{
  colback=csoft,colframe=cwarn,
  boxrule=0pt,leftrule=3pt,arc=2pt,
  fontupper=\small,
  enhanced,breakable,
  title=#1,
  coltitle=cprimary,fonttitle=\bfseries\small,
}
\newtcolorbox{calloutnote}[1][]{
  colback=csoft,colframe=cprimary,
  boxrule=0pt,leftrule=3pt,arc=2pt,
  fontupper=\small,
  enhanced,breakable,
  title=#1,
  coltitle=cprimary,fonttitle=\bfseries\small,
}

% --- Hyperliens (à charger en dernier ou presque) ---------------------
\usepackage[colorlinks=true,linkcolor=cprimary,citecolor=cprimary,
            urlcolor=cprimary]{hyperref}

% --- Titres de sections ----------------------------------------------
\usepackage{titlesec}
\titleformat{\section}{\Large\bfseries\color{cprimary}}{\thesection}{0.6em}{}
\titleformat{\subsection}{\large\bfseries\color{cprimary!85!black}}{\thesubsection}{0.5em}{}
\titleformat{\subsubsection}{\normalsize\bfseries\color{cprimary!75!black}}{\thesubsubsection}{0.4em}{}

% --- Code et symboles -------------------------------------------------
\usepackage{verbatim}
\usepackage{listings}
\lstset{basicstyle=\ttfamily\small,breaklines=true,frame=single,
        backgroundcolor=\color{csoft},rulecolor=\color{cline}}

% --- Notations actuarielles -----------------------------------------
\newcommand{\Espe}{\mathbb{E}}
\newcommand{\Prob}{\mathbb{P}}
\newcommand{\Mes}{\mathbb{Q}}
\newcommand{\indic}[1]{\mathbf{1}_{\{#1\}}}
\newcommand{\Norm}{\mathcal{N}}
\newcommand{\PD}{\mathrm{PD}}
\newcommand{\LGD}{\mathrm{LGD}}
\newcommand{\ETL}{\mathrm{ETL}}
\newcommand{\DF}{\mathrm{DF}}
\DeclareMathOperator{\Var}{Var}

% --- Macros figure et tableau ----------------------------------------
\makeatletter
% Strip ".png" from filename to build a clean label fig:<basename>
\newcommand{\fig}[3]{%
  \begin{figure}[H]
    \centering
    \includegraphics[width=#1\linewidth]{#2}
    \caption{#3}
    \label{fig:\@stripPng#2\@nil}
  \end{figure}}
\def\@stripPng#1.png\@nil{#1}
\makeatother

% =====================================================================
% TITRE
% =====================================================================
\title{
  \vspace{-1.2cm}
  {\color{cprimary}\Large Pricing et gestion des risques d'un CDO synthétique}\\[2pt]
  {\Large \textit{Pipeline de valorisation et d'analyse de bout en bout}}\\[6pt]
  {\normalsize Note de travail — Quant / Risk desk}
}
\author{
  \textsc{Tom Nguyen Duc Khanh Trinh}\\
  \small Master 2 Sciences Actuarielles, ISFA --- Université Claude Bernard Lyon 1\\
  \small Promotion 2025--2026
}
\date{26 avril 2026}

% =====================================================================
\begin{document}
\maketitle
\thispagestyle{empty}

\vspace{-0.3cm}
\begin{center}
\fbox{\parbox{0.94\textwidth}{\small
\textbf{Type de document :} note de travail modèle (desk de Crédit Structuré / Validation des Modèles).\\
\textbf{Public visé :} pricing desk, validation des modèles, gestion des risques.\\
\textbf{Références principales :} Mounfield (2009), Hull (2018, ch.~25), JP Morgan (2004), ISDA 2003.\\
\textbf{Pool de référence :} 651 émetteurs (CDS 5Y, équipondérés), date de valorisation 10/09/2021.\\
\textbf{Modèle retenu :} copule gaussienne à un facteur, $\rho = 0{,}30$.}}
\end{center}

% =====================================================================
\section*{Résumé exécutif}
\addcontentsline{toc}{section}{Résumé exécutif}

Cette note documente la mise en \oe uvre d'un \textbf{pipeline de pricing et de gestion
des risques de qualité dealer} pour les tranches d'un CDO synthétique. Le pipeline est
articulé en sept blocs : \emph{term-sheet $\to$ données de marché (CDS) $\to$ courbes de
crédit marginales (bootstrap ISDA) $\to$ modèle de dépendance (copule gaussienne à un
facteur) $\to$ flux de tranche (Monte Carlo) $\to$ calibration (corrélation de base)
$\to$ sensibilités et stress tests $\to$ tableau de bord portefeuille}.

Le résultat \textbf{passe l'ensemble des contrôles de validation} sur les quatre
critères critiques : repricing du bootstrap ISDA, cohérence LHP--MC, monotonicité de la
structure de capital, et raisonnement sur le signe de la corrélation. Les
\emph{primes équitables} obtenues (à $\rho = 30\%$, 5~ans, 10~000 trajectoires) :
Equity \SI{1944}{bps}, Mezzanine \SI{501}{bps}, Junior senior \SI{208}{bps}, Senior
\SI{89}{bps} --- monotonicité de la structure de capital confirmée.

\tableofcontents
\newpage

% =====================================================================
\section{Cadre théorique et données de marché}\label{sec:framework}

Le pipeline débute par la fixation du \textbf{term-sheet} de l'opération et la
collecte des \textbf{intrants de marché minimaux} : courbe d'actualisation, spreads
CDS aux maturités standards et hypothèse de récupération. Pour un CDO synthétique
standard de type iTraxx/CDX, le term-sheet comprend un portefeuille de référence
$\{i = 1, \dots, N\}$ de notionnels $N_i$ avec leurs cotations CDS, un calendrier de
coupons (trimestriel, base \emph{act/360}), une maturité $T$, une structure de capital
$[A_j, D_j]$ pour chaque tranche, ainsi qu'un coupon couru et/ou une prime upfront
(pour la tranche equity). La date de valorisation correspond au dernier instantané
du fichier \texttt{cds.csv}.

\subsection{Environnement de calcul}
Le \emph{notebook} ne dépend que de \texttt{numpy}, \texttt{pandas},
\texttt{matplotlib} et \texttt{seaborn}, sans recours à \texttt{scipy} pour assurer
la \textbf{portabilité} en environnement de validation. Les fonctions statistiques
standards (CDF et PPF de la loi normale) sont implémentées en interne via
\texttt{math.erf} et l'approximation d'Acklam (2003). La graine du générateur
pseudo-aléatoire est fixée à 42, garantissant la \textbf{reproductibilité} des
résultats --- exigence incontournable en validation modèle.

\subsection{Structure de capital --- term-sheet standard}
La structure de capital utilisée est une version compactée de l'iTraxx Europe Main
(tableau~\ref{tab:tranches}, figure~\ref{fig:fig01_capital_structure}).

\begin{table}[H]
\centering
\begin{tabular}{lccccp{6cm}}
\toprule
\textbf{Tranche} & $A$ & $D$ & \textbf{Largeur} & \textbf{Profil économique} \\
\midrule
Equity        & 0\,\%  & 3\,\%  & 3\,\% & Premier porteur de pertes --- \emph{long correlation} \\
Mezzanine     & 3\,\%  & 7\,\%  & 4\,\% & Short gamma --- sensible au \emph{default clustering} \\
Junior Senior & 7\,\%  & 10\,\% & 3\,\% & Double protection --- quasi rho-neutre \\
Senior        & 10\,\% & 15\,\% & 5\,\% & Protection de queue --- \emph{short correlation} \\
\bottomrule
\end{tabular}
\caption{Structure de capital iTraxx-like utilisée dans le pipeline.}
\label{tab:tranches}
\end{table}

\fig{0.85}{fig01_capital_structure.png}{Échelle de la structure de capital
d'un CDO synthétique standard (style iTraxx) : Equity 0--3\,\% (bleu), Mezz 3--7\,\%
(vert), Junior senior 7--10\,\% (orange), Senior 10--15\,\% (violet).}

La perte de tranche $[A,D]$ est définie selon la formule classique
\begin{equation}
L^{[A,D]}(t) \;=\; \min\!\bigl(\max(L^{\text{pool}}(t) - A,\,0),\; D - A\bigr),
\label{eq:tranche_loss}
\end{equation}
où la perte agrégée du \emph{pool} s'écrit
$L^{\text{pool}}(t) = \frac{1}{N}\sum_{i=1}^N \indic{\tau_i \le t}\cdot(1-R)$, avec
$\tau_i$ le temps de défaut de l'émetteur $i$. La tranche $[A,D]$ ne supporte de
pertes qu'entre les seuils $A$ (\emph{attachment}) et $D$ (\emph{detachment}) : en
deçà de $A$ elle reste intacte, au-delà de $D$ elle est entièrement absorbée.

\subsection{Données de marché --- CDS et pool de référence}
Le fichier \texttt{cds.csv} contient les CDS \emph{single-name} aux maturités
$\{1, 2, 3, 5, 7, 10\}$~ans pour plusieurs centaines d'émetteurs nord-américains et
européens. Les données sont \emph{pivotées} en format large
(\emph{Date $\times$ Ticker $\times$ Tenor}) et les spreads convertis de points de
base en \textbf{décimal} (division par 10\,000). Le pool de référence est extrait
des colonnes \texttt{PX5\_*} (CDS 5Y), la maturité de l'opération étant fixée par
défaut à $T = 5$~ans, conformément à la pratique iTraxx Main.

\begin{table}[H]
\centering
\begin{tabular}{lr}
\toprule
\textbf{Indicateur} & \textbf{Valeur} \\
\midrule
Trame CDS                          & 1\,747 dates $\times$ 6\,610 colonnes \\
Pool 5Y de référence (initial)     & 661 émetteurs \\
Pool 5Y après bootstrap            & 651 émetteurs valides \\
Date de valorisation               & 10/09/2021 \\
\bottomrule
\end{tabular}
\caption{Volumétrie des données de marché.}
\label{tab:market_data}
\end{table}

\subsection{Courbe d'actualisation et hypothèse de récupération}
On utilise un \textbf{taux sans risque continu et constant} $r = 2\%$ (substituable
par une courbe OIS/SOFR réelle en production) :
\begin{equation}
\DF(t) = e^{-r\,t}.
\end{equation}
Le taux de récupération est fixé à $R = 40\%$ (convention ISDA pour la dette senior
non garantie), soit une perte en cas de défaut $\LGD = 1 - R = 60\%$. Au § VI, $R$
sera choqué à 25\,\% puis 10\,\% pour évaluer l'impact d'une récupération
\emph{distressed}.

% =====================================================================
\section{Courbes de crédit marginales}\label{sec:credit_curves}

Chaque émetteur de référence requiert une \textbf{courbe de survie} $Q_i(t)$
fournissant la probabilité de non-défaut jusqu'à $t$ sous la mesure risque-neutre
$\Mes$. La PD à l'horizon $T$ s'en déduit immédiatement : $p_i(T) = 1 - Q_i(T)$.
Deux constructions sont possibles : (i) raccourci à \emph{hasard plat} (rapide,
ignore la structure par terme) et (ii) bootstrap par morceaux constants
\emph{ISDA-compatible} (standard de production).

\subsection{Approximation à hasard plat (référence pédagogique)}
On applique la formule du \emph{credit triangle} de Hull,
\begin{equation}
\lambda_i \approx \frac{s_i}{1 - R}.
\label{eq:credit_triangle}
\end{equation}
Cette formule est exacte dans la limite d'une fréquence de coupons infinie et d'un
taux sans risque nul. On l'utilise comme référence de comparaison.

\begin{table}[H]
\centering
\begin{tabular}{lc}
\toprule
Pool $N$                       & 651 (équipondéré) \\
PD~5Y plate (moyenne)          & 5,574\,\% \\
PD~5Y plate (étendue)          & [0,86\,\% ; 47,78\,\%] \\
\bottomrule
\end{tabular}
\caption{Statistiques sommaires du pool sous hasard plat.}
\end{table}

\subsection{Bootstrap ISDA par morceaux constants}
Soit $\lambda_i(t)$ une fonction constante par morceaux sur les n\oe uds
$0 = T_0 < T_1 < \cdots < T_K$ correspondant aux maturités $\{1, 2, 3, 5, 7, 10\}$~ans.
La probabilité de survie admet la forme fermée
\begin{equation}
Q_i(t) = \exp\!\left(-\sum_{k=1}^K \lambda_i^{(k)}\,\Delta_k(t)\right),
\quad
\Delta_k(t) = \max\!\bigl(\min(T_k, t) - T_{k-1},\, 0\bigr).
\label{eq:Q_pwc}
\end{equation}
Pour chaque maturité $T_k$, on résout par dichotomie l'équation traduisant la
nullité de la NPV du contrat CDS :
\begin{equation}
\underbrace{\int_0^{T_k} s\,\DF(u)\,Q(u)\,du}_{\text{jambe de prime}}
\;=\;
\underbrace{(1-R)\int_0^{T_k} \DF(u)\bigl(-dQ(u)\bigr)}_{\text{jambe de protection}}.
\label{eq:CDS_bootstrap}
\end{equation}
L'algorithme itère sur $k = 1, \dots, K$ : à chaque étape on fixe les hasards
précédents $\lambda^{(1)}, \dots, \lambda^{(k-1)}$ et l'on résout pour
$\lambda^{(k)}$ --- ce qui correspond exactement à l'\emph{ISDA Standard Model}
mis en \oe uvre dans les systèmes \emph{dealer}.

\begin{table}[H]
\centering
\begin{tabular}{lr}
\toprule
Émetteurs bootstrappés        & 651 (\textasciitilde 30~s wall-clock) \\
PD~5Y bootstrap (moyenne)     & 5,717\,\% \\
PD~5Y bootstrap (écart-type)  & 5,105\,\% \\
\bottomrule
\end{tabular}
\caption{Statistiques après bootstrap ISDA par morceaux constants.}
\end{table}

\subsection{Auto-cohérence et structure par terme}

Après bootstrap, deux contrôles doivent passer : (i) l'\emph{erreur de repricing}
--- avec la courbe construite, la NPV du CDS 5Y d'entrée doit satisfaire
$|\text{NPV}| < 10^{-5}$ (bruit numérique) ; (ii) la \emph{forme de la structure
par terme} par bucket de notation --- les émetteurs bien notés (AAA/AA) doivent
présenter une courbe basse et plate, tandis que les émetteurs faiblement notés
(CCC/D) doivent montrer une courbe fortement croissante traduisant
l'\emph{anticipation de détérioration de crédit}.

\begin{callout}[Résultat de la validation du bootstrap]
N tested = 300 émetteurs ; max $|\text{NPV}| = \SI{8.33e-17}{}$ ; mean
$|\text{NPV}| = \SI{2.95e-17}{}$ ; \textbf{100\,\% de réussite} sur le critère
$< 10^{-5}$. Le bootstrap est confirmé \textbf{ISDA-compatible}.
\end{callout}

\fig{0.85}{fig02_bootstrap_repricing.png}{Histogramme de l'erreur de repricing
des CDS~5Y après bootstrap (unité : $\times 10^{-6}$ du notionnel). Toutes les
valeurs sont extrêmement proches de zéro.}

\fig{0.95}{fig03_hazard_termstruct.png}{À gauche : structure par terme moyenne du
hasard par bucket de notation --- les bonnes notations sont basses et plates ; les
faibles sont à pente positive marquée. À droite : distribution des PD~5Y sur
l'ensemble du pool, moyenne \textasciitilde 5,7\,\%, étendue [0,86\,\% ; 47,78\,\%].}

% =====================================================================
\section{Modèle de dépendance --- copule gaussienne à un facteur}\label{sec:copula}

Le \textbf{module de dépendance} est le c\oe ur du pricing CDO : la prime de
tranche dépend de la \emph{distribution jointe des pertes} sur l'ensemble du pool,
non des seules PD marginales. Suivant Li (2000) et Vasicek (1987), on associe à
chaque émetteur un rendement d'actif latent
\begin{equation}
A_i = \sqrt{\rho}\,M + \sqrt{1-\rho}\,\varepsilon_i,
\qquad M, \varepsilon_i \sim \Norm(0,1) \text{ i.i.d.}
\label{eq:latent_return}
\end{equation}
Le défaut survient lorsque $A_i < \Phi^{-1}(p_i(T))$. Conditionnellement au facteur
commun $M = m$, la PD individuelle s'écrit
\begin{equation}
p_i(T \mid M = m) = \Phi\!\left(
\frac{\Phi^{-1}(p_i(T)) - \sqrt{\rho}\,m}{\sqrt{1-\rho}}\right).
\label{eq:cond_pd}
\end{equation}
Intuition : lorsque $M$ se dégrade ($m$ très négatif), tous les $p_i(T \mid M)$
sautent simultanément, générant le \emph{default clustering} qui produit la
queue gauche épaisse de la distribution des pertes.

\begin{table}[H]
\centering
\begin{tabular}{rc}
\toprule
$p_i(T \mid M = -2)$ & 25,57\,\% \quad (régime de stress, défauts groupés) \\
$p_i(T \mid M = 0)$  & 2,47\,\% \quad ($\approx$ PD marginale recalibrée) \\
$p_i(T \mid M = +2)$ & 0,05\,\% \quad (régime favorable, quasi-zéro défaut) \\
\bottomrule
\end{tabular}
\caption{PD conditionnelles, calculées avec $p_{\text{marg}} = 5\%$ et $\rho = 0{,}30$.}
\end{table}

\subsection{Limite \emph{Large Homogeneous Pool} (Vasicek 1987)}
Sous l'hypothèse d'un pool homogène ($p_i \equiv p$) et $N \to \infty$, la loi des
grands nombres entraîne
\begin{equation}
L^{\text{pool}}(T) \;\xrightarrow{p}\;
\LGD \cdot \Phi\!\left(
\frac{\Phi^{-1}(p) - \sqrt{\rho}\,M}{\sqrt{1-\rho}}\right).
\label{eq:LHP_loss}
\end{equation}
On en déduit la \emph{fonction de répartition} de la perte du pool sous LHP :
\begin{equation}
F_L(\ell) = \Phi\!\left(
\frac{\sqrt{1-\rho}\,\Phi^{-1}(\ell/\LGD) - \Phi^{-1}(p)}{\sqrt{\rho}}\right).
\label{eq:LHP_cdf}
\end{equation}
C'est la formule de Vasicek (1987), également au fondement de la formule de poids
de risque IRB de Bâle~II.

\subsection{Corrélation et risque de queue de tranche}

\fig{0.95}{fig04_loss_distrib_corr.png}{Comparaison de la distribution de pertes du
pool sous deux régimes de corrélation. À gauche $\rho = 0{,}10$ : distribution
serrée, senior protégée. À droite $\rho = 0{,}40$ : queue droite épaissie, senior
exposée au risque de queue. Pointillé rouge = $\Espe[L]$ ; zones colorées = tranches
de la structure de capital.}

À $\rho = 0{,}10$, la distribution se concentre autour de $\Espe[L] = p\cdot\LGD
\approx 3\%$ avec une queue mince --- la senior est en sécurité. À $\rho = 0{,}40$,
la queue droite s'épaissit considérablement et $\Prob(L > 10\%)$ croît d'un ordre
de grandeur. C'est précisément la raison pour laquelle la \textbf{tranche senior
est short corrélation} et la \textbf{tranche equity est long corrélation},
propriété que confirmera le signe de Rho01 au § VI.

% =====================================================================
\section{Pricing par Monte Carlo}\label{sec:pricing}

Disposant des PD marginales $\{p_i(T)\}$ et de la corrélation $\rho$, le pricing
d'une tranche $[A, D]$ procède en trois étapes :

\begin{enumerate}
\item Simulation des \textbf{temps de défaut} $\tau_i$ par copule gaussienne et
      marginales exponentielles ;
\item Calcul de la \textbf{perte attendue de tranche}
      $\ETL_k = \Espe[L^{[A,D]}(t_k)]$ sur la grille de coupons ;
\item Résolution du \textbf{spread équitable courant} $s^\star$ tel que
      PV(protection) = PV(prime) :
\begin{equation}
s^\star =
\frac{\displaystyle\sum_k \DF(t_k)\,(\ETL_k - \ETL_{k-1})}
     {\displaystyle\sum_k \DF(t_k)\,\Delta t_k\,
       \bigl((D-A) - \tfrac{1}{2}(\ETL_k + \ETL_{k-1})\bigr)}.
\label{eq:fair_spread}
\end{equation}
Le numérateur est la jambe de protection ; le dénominateur est l'annuité (DV01) de
la tranche avec notionnel restant approximé par la règle des trapèzes.
\end{enumerate}

\subsection{Simulateur de temps de défaut et grille de paiements}
La copule gaussienne génère $Y_i = \sqrt{\rho}\,M + \sqrt{1-\rho}\,\varepsilon_i$,
puis $U_i = \Phi(Y_i)$, et enfin $\tau_i = -\log(1 - U_i)/\lambda_i$ (marginale
exponentielle). Le même \emph{tau\_world} est réutilisé pour toutes les tranches
afin d'éliminer le bruit MC entre tranches. La grille de paiements est trimestrielle,
$\Delta t = 0{,}25$~an, $T = 5$~ans, soit 20 coupons.

\begin{table}[H]
\centering
\small
\begin{tabular}{lcccrr}
\toprule
\textbf{Tranche} & $[A, D]$ & \textbf{Largeur} &
\textbf{$\ETL_{5Y}$ (notional)} & \textbf{Spread équitable (bps)} \\
\midrule
Equity        & 0--3\,\%   & 3\,\% & 0,01834 & \textbf{1\,943,6} \\
Mezz          & 3--7\,\%   & 4\,\% & 0,00909 & \textbf{500,6} \\
Junior sen.   & 7--10\,\%  & 3\,\% & 0,00304 & \textbf{207,8} \\
Senior        & 10--15\,\% & 5\,\% & 0,00222 & \textbf{88,9} \\
\bottomrule
\end{tabular}
\caption{Résultats de pricing à $\rho = 0{,}30$, $T = 5$~ans, 10\,000 trajectoires
MC, graine 42.}
\label{tab:pricing}
\end{table}

\fig{0.95}{fig05_etl_paths_spreads.png}{À gauche : trajectoires d'$\ETL$ des
quatre tranches sur la grille trimestrielle (5~ans) --- la tranche Equity accumule
les pertes le plus rapidement, la senior ne « s'éveille » qu'au dernier trimestre.
À droite : spread équitable (bps) --- décroissance forte avec la séniorité
(1\,944~$\to$~89~bps).}

\begin{callout}[Lecture du tableau]
La tranche Equity « consomme » plus de 60\,\% de son propre notionnel en perte
attendue --- le spread équitable de 1\,944~bps reflète exactement sa position de
\emph{first-loss}. La senior n'absorbe que 4,4\,\% de son notionnel propre,
spread 89~bps. La \textbf{monotonicité de la structure de capital} est confirmée :
Equity~$\gg$~Mezz~$>$~Junior sen.~$>$~Senior.
\end{callout}

\subsection{Réduction de variance par variables antithétiques}
Pour chaque tirage du facteur commun $m$, on ajoute le tirage $-m$ : la symétrie
des actifs latents conduit à des estimateurs sans biais et de variance réduite.
Le facteur de réduction de variance s'écrit
$\text{VRF} = \sigma^2_{\text{crude}} / \sigma^2_{\text{antithetic}}$. Sur le
pool ici considéré ($N = 651$, $10^4$ trajectoires), la VRF observée est
$\approx 0{,}98 \times$ : la technique apporte peu lorsque le pool est suffisamment
grand. Elle reste utile pour les pools réduits ou les corrélations élevées.

\fig{0.7}{fig06_antithetic_var.png}{Erreur standard de l'estimateur MC pour la
tranche Mezzanine [3\,\%; 7\,\%], comparaison \emph{crude} vs \emph{antithetic}
selon le nombre de trajectoires. Pente $\approx -1/2$ en échelle log-log,
correspondant au taux de convergence $\sqrt{n}$.}

\subsection{Référence semi-analytique LHP --- benchmarking}
La forme close LHP fournit une référence indépendante du moteur MC. Le notebook
évalue LHP au PD moyen avec $n_{gh} = 64$ (Gauss--Hermite) et 96~points de Legendre.

\begin{table}[H]
\centering
\begin{tabular}{lcrrr}
\toprule
\textbf{Tranche} & $[A, D]$ & \textbf{$\ETL$ LHP} & \textbf{$\ETL$ MC} & \textbf{Diff (bps)} \\
\midrule
Equity      & 0--3\,\%   & 0,01847 & 0,01834 & $+1{,}3$ \\
Mezz        & 3--7\,\%   & 0,01403 & 0,00909 & $+49{,}4$ \\
Junior sen. & 7--10\,\%  & 0,00746 & 0,00304 & $+44{,}2$ \\
Senior      & 10--15\,\% & 0,00911 & 0,00222 & $+68{,}8$ \\
\bottomrule
\end{tabular}
\caption{Comparaison LHP semi-analytique vs Monte Carlo.}
\end{table}

\fig{0.75}{fig07_lhp_vs_mc.png}{Balayage en corrélation $\rho \in [0{,}02 ; 0{,}90]$
pour la tranche Equity [0--3\,\%] : LHP semi-analytique (trait plein) et Monte Carlo
5\,000~trajectoires (points). Les deux méthodes coïncident à l'intérieur du bruit MC
jusqu'à $\rho \le 0{,}6$.}

\begin{calloutwarn}[Observation cruciale]
La tranche Equity dépend quasi exclusivement de la PD moyenne du pool ; LHP et MC
coïncident (erreur $< 5\,\%$). En revanche, Mezz et Senior sont sensibles à
l'\emph{hétérogénéité}, et le pool réel présente $p_i$ étalés sur
[0,86\,\% ; 47,78\,\%]. Le MC capture une queue plus épaisse que la LHP homogène
--- déviations de 10 à 70\,\% (Mounfield §6.5, O'Kane 2008). Remèdes en production :
\textbf{LHP par bucket} ou \textbf{récursion de Hull--White}, ou MC seul pour
Mezz/Senior.
\end{calloutwarn}

% =====================================================================
\section{Calibration --- corrélation de base}\label{sec:calib}

Le marché cote des \emph{spreads de tranches d'indice}. Pour reproduire les
cotations, le desk ne cote pas $\rho$ directement, mais l'extrait par inversion
du modèle. Deux conventions coexistent :

\begin{itemize}
\item \textbf{Corrélation composée} $\rho^{\text{comp}}_j$ : pour chaque tranche
      $[A_j, D_j]$, on résout $\rho$ tel que le spread modèle = spread marché.
      Inconvénient : non-unicité, smile dégradé pour la mezzanine.
\item \textbf{Corrélation de base} $\rho^B(D)$ (JP Morgan 2004) : pour chaque
      $D$ standard, on résout $\rho$ tel que l'$\ETL$ de l'\emph{equity agrégé}
      $[0, D]$ corresponde à la cible. Avantages : monotonie en $D$, interpolation
      possible, standard de l'industrie.
\end{itemize}

Une tranche \emph{bespoke} $[A, D]$ se valorise alors par
\begin{equation}
\ETL^{[A, D]} = \ETL^{[0, D]}\bigl(\rho^B(D)\bigr) - \ETL^{[0, A]}\bigl(\rho^B(A)\bigr).
\label{eq:bespoke_etl}
\end{equation}

\subsection{Calibration de la courbe BC (round-trip pédagogique)}

Le notebook simule un \emph{smile} typique iTraxx (BC croissante avec $D$,
reflétant la prime de risque systémique) :
\begin{equation}
\rho^B(D) :\;\{3\% \to 20\%,\; 7\% \to 30\%,\; 10\% \to 37\%,\; 15\% \to 45\%\}.
\end{equation}
Pour chaque $D$, l'$\ETL$ cible est généré par LHP-au-quote ; la BC est ensuite
récupérée par dichotomie sur le même objectif LHP. C'est la procédure de
validation du calibrateur. En production, la cible $\ETL$ est dérivée du spread
marché coté via la formule de spread équitable du § IV.

\begin{calloutnote}[Correctif 04/2026 --- round-trip parfait]
La version précédente de \texttt{lhp\_etl} souffrait de \textbf{deux bugs
combinés} :
(a) mauvaise convention de Gauss--Hermite (utilisation de
\texttt{nodes*sqrt(2), wts/sqrt(pi)} adaptée au poids \emph{physiciste} $e^{-x^2}$
alors que \texttt{hermegauss} de \texttt{hermite\_e} utilise le poids
\emph{probabiliste} $e^{-x^2/2}$ --- gonflant la variance du facteur latent d'un
facteur $\sqrt{2}$) ;
(b) formule interne erronée --- application de
\texttt{lhp\_cond\_cdf}$(\ell, p_c(m), \rho)$ dans la boucle Hermite,
\emph{double-comptant} le facteur commun $M$.
Après remplacement par la formule de Vasicek standard
$L \mid m = \LGD \cdot p_c(m)$ et une dichotomie indépendante du sens de
monotonie, le \textbf{round-trip BC restitue exactement les cotations d'entrée}
(20\,\% / 30\,\% / 37\,\% / 45\,\%) à la précision $10^{-7}$.
\end{calloutnote}

\begin{table}[H]
\centering
\begin{tabular}{lrrr}
\toprule
\textbf{Tranche de base} & \textbf{$\ETL$ cible (bps)} & \textbf{BC marché (entrée)} & \textbf{BC restituée} \\
\midrule
0--3\,\%   & 203,9 & 20,00\,\% & 20,00\,\% \\
0--7\,\%   & 267,6 & 30,00\,\% & 30,00\,\% \\
0--10\,\%  & 280,6 & 37,00\,\% & 37,00\,\% \\
0--15\,\%  & 300,8 & 45,00\,\% & 45,00\,\% \\
\bottomrule
\end{tabular}
\caption{Round-trip de la calibration de corrélation de base après correction.}
\end{table}

\fig{0.75}{fig08_bc_curve.png}{Courbe de corrélation de base $\rho^B(D)$ après
correction --- la BC restituée (en vert) co\"incide avec les cotations de marché
(en pointillé bleu) à 20\,\%/30\,\%/37\,\%/45\,\%, validant le calibrateur à la
précision $10^{-7}$.}

\subsection{Surface de spread équitable pour tranches bespoke}
Une fois $\rho^B(D)$ établi, on interpole linéairement par morceaux pour valoriser
toute tranche bespoke $[A, D]$. La sortie est une \textbf{surface} sur la grille
$(A,\text{largeur})$ --- outil de pricing rapide pour le desk.

\begin{table}[H]
\centering
\begin{tabular}{lrr}
\toprule
\textbf{Tranche bespoke} & \textbf{$\ETL$ (bps)} & \textbf{Spread (bps)} \\
\midrule
0--5\,\%   & 244,9 & 1\,060,1 \\
4--8\,\%   & 44,7  & 242,1   \\
6--12\,\%  & 26,7  & 96,5    \\
10--15\,\% & 25,4  & 109,8   \\
15--25\,\% & 25,6  & 55,4    \\
\bottomrule
\end{tabular}
\caption{Échantillon de cotations bespoke après correction du calibrateur BC.}
\end{table}

\fig{0.95}{fig09_bespoke_surface.png}{Surface de spread équitable (bps) sur la
grille \emph{(attachment $\times$ largeur)} après correction de \texttt{lhp\_etl}.
Rouge = spread élevé (tranches basses, étroites) ; vert = spread faible (senior,
larges). Étoiles blanches = tranches standard de style iTraxx.}

% =====================================================================
\section{Analyse des risques --- Greeks et stress tests}\label{sec:risk}

Après pricing, la gestion des risques se concentre sur la \textbf{mesure des
sensibilités} et le \textbf{stress test scénarisé}. L'ensemble minimal de Greeks
de tranche est :

\begin{table}[H]
\centering
\begin{tabular}{lll}
\toprule
\textbf{Indicateur} & \textbf{Définition} & \textbf{Instrument de couverture} \\
\midrule
CS01  & $\partial s^\star / \partial(\text{spread}+1\text{bp})$ & CDS \emph{single-name} (parallèle +1bp) \\
Rho01 & $\partial s^\star / \partial \rho$              & Correlation swap / Index vs tranche \\
DV01  & $\partial s^\star / \partial r$                 & IRS (swap de taux) \\
JtD   & $\Delta s^\star$ si 1 émetteur défaute immédiatement & Couverture LGD \\
\bottomrule
\end{tabular}
\caption{Greeks de tranche.}
\end{table}

\subsection{Greeks au niveau tranche}
La méthode \emph{bump-and-reprice} est employée avec la \textbf{même graine} pour
éliminer le bruit MC entre baseline et scénario choqué.

\begin{table}[H]
\centering
\small
\begin{tabular}{lrrrrr}
\toprule
\textbf{Tranche} & \textbf{Spread (bps)} & \textbf{CS01} & \textbf{Rho01} & \textbf{DV01} & \textbf{JtD} \\
\midrule
Equity      & 1\,952,6 & $+21{,}93$ & $-37{,}50$ & $+0{,}03$  & $+145{,}5$ \\
Mezz        & 503,3   & $+9{,}64$  & $-0{,}73$  & $-0{,}02$  & $+12{,}9$  \\
Junior sen. & 202,7   & $+5{,}63$  & $+3{,}90$  & $-0{,}01$  & $+3{,}9$   \\
Senior      & 85,4    & $+3{,}02$  & $+3{,}83$  & $-0{,}01$  & $+1{,}5$   \\
\bottomrule
\end{tabular}
\caption{Greeks par tranche. CS01 et DV01 en bps/+1bp ; Rho01 en bps/+1\% ; JtD
en bps par défaut.}
\label{tab:greeks}
\end{table}

\fig{0.95}{fig10_greeks.png}{Quatre Greeks par tranche (CS01 / Rho01 / DV01 / JtD).
Equity domine CS01 et JtD ; Rho01 change de signe entre Equity (négatif,
\emph{long correlation}) et Senior (positif, \emph{short correlation}) --- la
signature classique.}

\begin{callout}[Lecture par logique économique]
\textbf{Equity} présente le plus grand CS01 ($+21{,}9$~bps) et un Rho01 fortement
\textbf{négatif} ($-37{,}5$~bps) : l'equity bénéficie d'une hausse de corrélation
car $\Prob(L \ge 3\%)$ baisse plus vite que ne croît la perte attendue
($\Rightarrow$ \emph{long correlation}). \textbf{Senior} affiche un Rho01
\textbf{positif} ($+3{,}8$~bps) : la queue s'épaissit avec $\rho$, dégradant la
senior ($\Rightarrow$ \emph{short correlation}). \textbf{JtD} : impact massif sur
l'equity ($+145{,}5$~bps pour 1 défaut) et marginal sur la senior ($+1{,}5$~bps),
celle-ci étant « profondément hors-tranche » d'un défaut isolé. \textbf{DV01}
reste faible sur toutes les tranches ($\le 0{,}03$~bps) --- les tranches ne sont
pas un instrument de taux.
\end{callout}

\subsection{CS01 par émetteur et concentration de la couverture}
Sous l'approximation \emph{mean-field} LHP, le CS01 d'un émetteur $k$ vaut
\begin{equation}
\text{CS01}_k =
\frac{\partial \ETL}{\partial p_{\text{avg}}}
\cdot
\frac{\partial p_k}{\partial s_k}
\cdot
\frac{1}{N}.
\label{eq:cs01_namelevel}
\end{equation}
Les contributeurs principaux sont les émetteurs au spread le plus élevé. En
configuration \emph{mean-field} (poids $1/N$), le top-10 représente \textasciitilde
1,6\,\% et le top-50 \textasciitilde 8,0\,\% du CS01 du pool --- répartition
quasi-uniforme par construction. Avec un gradient hétérogène (production), le
top-10 capte typiquement 30--40\,\% du CS01 equity et le desk ne couvre que
\textasciitilde 10 CDS pour aplatir la tranche equity.

\begin{table}[H]
\centering
\small
\begin{tabular}{lrr}
\toprule
\textbf{Ticker} & \textbf{PD~5Y} & \textbf{CS01 Equity} ($\times 10^{-4}$) \\
\midrule
NOVNVX  & 0,876\,\% & 3,95 \\
AGSBB   & 0,952\,\% & 3,95 \\
XL      & 1,075\,\% & 3,95 \\
TOYOTA  & 1,325\,\% & 3,94 \\
NESNVX  & 1,418\,\% & 3,93 \\
DBHNGR  & 1,470\,\% & 3,93 \\
HENKEL  & 1,475\,\% & 3,93 \\
UNANA   & 1,497\,\% & 3,93 \\
NTT     & 1,512\,\% & 3,93 \\
EJRAIL  & 1,545\,\% & 3,93 \\
\bottomrule
\end{tabular}
\caption{Top-10 contributeurs CS01 equity (PD bootstrappées).}
\end{table}

\fig{0.95}{fig11_hedge_concentration.png}{À gauche : top-20 contributeurs CS01
pour la tranche Equity (LHP \emph{mean-field}). À droite : CS01 vs PD~5Y par
émetteur, par tranche --- l'equity domine en magnitude, la senior n'ajuste
qu'à la marge.}

\subsection{Matrice de stress tests}
La grille de chocs est \textbf{déterministe} (sans probabilité d'occurrence ;
conformément aux exigences Bâle~III et Fed CCAR) :

\begin{table}[H]
\centering
\begin{tabular}{lll}
\toprule
\textbf{Scénario} & \textbf{Choc} & \textbf{Interprétation} \\
\midrule
Spread +50 / +100 / +200 bp     & Décalage de hasard parallèle      & Retournement de cycle \\
Recovery 40\,\%~$\to$~25\,\% / 10\,\% & Récupération en détresse          & Hausse brutale du LGD \\
$\rho$ +20pt / +50pt            & Saut de corrélation               & Default clustering --- événement systémique \\
Crise type 2008                 & Combiné : spread +100bp, $R = 20\%$, $\rho$ +35pt & Triple choc à la GFC \\
\bottomrule
\end{tabular}
\caption{Grille de scénarios de stress.}
\end{table}

\fig{0.95}{fig12_stress_heatmap.png}{Carte de chaleur du choc de spread
$\Delta s^\star$ (bps vs base) sur 8 scénarios $\times$ 4 tranches.}

\begin{calloutwarn}[Règle empirique]
L'equity est la plus sensible aux chocs de spread et de recovery ; la senior est
la plus sensible aux chocs de corrélation (événement \emph{long correlation}). La
heatmap permet au risk manager d'identifier le scénario le plus douloureux par
tranche et de dimensionner la couverture en conséquence.
\end{calloutwarn}

% =====================================================================
\section{Analyse de portefeuille et tableau de bord}\label{sec:portfolio}

La dernière partie présente l'\textbf{analyse au niveau portefeuille} sur série
temporelle. Le LHP semi-analytique permet de constituer des instantanés mensuels
sur plus de 6~ans à coût numérique réduit, donnant accès à : l'évolution de la PD
moyenne du pool sur le cycle 2015--2021 ; le choc COVID-19 du T1~2020 et son
levier par tranche ; le \emph{regime shift} pré-crise (spread bas et stable)
$\to$ pic de mars 2020 $\to$ retour à la normale.

\subsection{$\ETL$ historiques par tranche (2015--2021)}

\begin{table}[H]
\centering
\begin{tabular}{lr}
\toprule
\textbf{Indicateur} & \textbf{Valeur} \\
\midrule
Instantanés mensuels         & 81 (de 01/2015 à 09/2021) \\
PD~5Y du pool, étendue       & 5,56\,\% -- 11,84\,\% \\
Date du pic de PD            & 01/04/2020 ($p_{\text{avg}} = 11,84\%$) \\
\bottomrule
\end{tabular}
\caption{Volumétrie de la série historique.}
\end{table}

\begin{table}[H]
\centering
\begin{tabular}{lcl}
\toprule
\textbf{Tranche} & \textbf{Levier COVID} & \textbf{Interprétation} \\
\midrule
Equity      & $1{,}21\times$ & Hausse modeste --- equity déjà près du fond \\
Mezz        & $1{,}33\times$ & Convexité émerge \\
Junior sen. & $1{,}39\times$ & Convexité nette \\
Senior      & $1{,}45\times$ & Plus forte hausse --- queue de risque ranimée \\
\bottomrule
\end{tabular}
\caption{Levier COVID = pic d'$\ETL$ / moyenne pré-COVID.}
\end{table}

\fig{0.95}{fig13_historical_etl.png}{Série temporelle 2015--2021. Panneau~1 : PD~5Y
moyenne du pool (zone rouge = choc COVID). Panneau~2 : $\ETL$ Equity et Mezz.
Panneau~3 : $\ETL$ Junior senior et Senior. La COVID frappe toutes les tranches
mais le levier relatif maximal s'observe sur la senior ($1{,}45\times$).}

\begin{calloutnote}[Insight]
Contre-intuitivement, la variation \emph{relative} de la senior est la plus forte
car celle-ci passe d'un état \emph{deep OTM} à proche-de-l'argent dès lors que le
risque de queue devient un scénario central. C'est une manifestation de la
\textbf{convexité} de la distribution de pertes et la raison pour laquelle la
senior présente un Rho01 nettement positif.
\end{calloutnote}

% =====================================================================
\section{Synthèse et liste de validation}\label{sec:summary}

\begin{table}[H]
\centering
\small
\begin{tabular}{lrrrr}
\toprule
\textbf{Indicateur} & \textbf{Equity} & \textbf{Mezz} & \textbf{Junior sen.} & \textbf{Senior} \\
\midrule
$A$--$D$                            & 0--3\,\% & 3--7\,\% & 7--10\,\% & 10--15\,\% \\
$\ETL_{5Y}$ (MC, du notionnel)      & 0,01834   & 0,00909   & 0,00304   & 0,00222   \\
$\ETL_{5Y}$ (LHP)                   & 0,01847   & 0,01403   & 0,00746   & 0,00911   \\
\textbf{Spread équitable (bps)}     & \textbf{1\,952,6} & \textbf{503,3} & \textbf{202,7} & \textbf{85,4} \\
CS01 (bps/+1bp)                     & $+21{,}93$ & $+9{,}64$ & $+5{,}63$ & $+3{,}02$ \\
Rho01 (bps/+1\% $\rho$)             & $-37{,}50$ & $-0{,}73$ & $+3{,}90$ & $+3{,}83$ \\
DV01 (bps/+1bp taux)                & $+0{,}03$  & $-0{,}02$ & $-0{,}01$ & $-0{,}01$ \\
JtD (bps/1 défaut)                  & $+145{,}5$ & $+12{,}9$ & $+3{,}9$  & $+1{,}5$  \\
\bottomrule
\end{tabular}
\caption{Rapport de risque final --- date de valorisation 10/09/2021,
$\rho = 0{,}30$, $T = 5$~ans, 651 émetteurs équipondérés.}
\end{table}

\subsection*{Liste de contrôle de validation}

\begin{itemize}
\item[\checkmark] \textbf{Repricing du bootstrap ISDA $< 10^{-5}$} --- 100,0\,\% des
émetteurs (max $|\text{NPV}| \approx \SI{8.3e-17}{}$).
\item[\checkmark] \textbf{Erreur relative LHP vs MC sur l'equity} = 0,7\,\%
(cible $< 20\,\%$).
\item[\checkmark] \textbf{Monotonicité de la structure de capital} : Equity~$>$~Mezz
~$>$~Junior sen.~$>$~Senior.
\item[\checkmark] \textbf{Cohérence des signes de corrélation} : Equity \emph{long
rho} (Rho01 $< 0$), Senior \emph{short rho} (Rho01 $> 0$).
\item[\checkmark] \textbf{Round-trip BC après correctif 04/2026} --- restitution
exacte des cotations marché à $10^{-7}$.
\end{itemize}

\subsection*{Notes et recommandations pour la production}

\begin{enumerate}
\item Pour les tranches Mezz/Senior, le LHP \emph{homogène} est attendu en
écart de 10 à 70\,\% par rapport au MC sur pool hétérogène. Préférer le MC ou
le \emph{bucketed-LHP} en marquage de production au-delà de l'equity.
\item La courbe BC est ici un round-trip pédagogique à partir de cotations
stylisées ; en production, alimenter le calibrateur avec des cotations
iTraxx/CDX réelles et utiliser une cible-spread (plutôt qu'une cible-$\ETL$).
\item Remplacer le taux constant $r = 2\%$ par une courbe OIS/SOFR réelle ;
substituer une recovery par notation à la valeur fixe de 40\,\%.
\item La graine du RNG est fixée à 42 pour la reproductibilité ; en production,
la graine et le nombre de trajectoires doivent être consignés dans l'\emph{audit
log}.
\item Le LHP en série temporelle peut servir au calcul de l'\emph{attribution
P\&L} via les CS01 et Rho01 quotidiens et s'intégrer au système de risque
front-office.
\end{enumerate}

% =====================================================================
\section*{Références}
\addcontentsline{toc}{section}{Références}

\begin{itemize}
\item Mounfield, C.~C. (2009). \emph{Synthetic CDOs: Modelling, Valuation and
      Risk Management}. Cambridge University Press.
\item Hull, J.~C. (2018). \emph{Options, Futures and Other Derivatives},
      9\textsuperscript{e} édition, chapitre 25. Pearson.
\item JPMorgan (2004). \emph{Credit Derivatives --- Base Correlation}. Note de
      recherche quantitative.
\item Vasicek, O. (1987). \emph{Probability of Loss on Loan Portfolio}. KMV
      Corporation.
\item Li, D.~X. (2000). On default correlation: A copula function approach.
      \emph{Journal of Fixed Income}, 9(4), 43--54.
\item O'Kane, D. (2008). \emph{Modelling Single-name and Multi-name Credit
      Derivatives}. Wiley Finance.
\item ISDA (2003). \emph{Credit Derivatives Definitions} ; \emph{ISDA Standard
      CDS Model documentation}.
\end{itemize}

\vfill
\begin{center}\small
\textit{Document compilé à partir du notebook \texttt{CDO\_Pipeline\_Pro.ipynb}
--- instantané du 10/09/2021, pool de 651 émetteurs, copule gaussienne à un
facteur ($\rho = 0{,}30$). Correctif appliqué le 26/04/2026 sur la fonction
\texttt{lhp\_etl} et le calibrateur de corrélation de base.}
\end{center}

\end{document}
